\section{Evaluierung}
\label{sec:evaluierung}

Der ausgewertete Lauf plant eine zehntägige Reise ab Stuttgart mit drei
Zielen und Unterkünften. Nach 600 Sekunden endet die Suche in Generation~464
von höchstens 500. Die vier Empfehlungen besuchen London, Barcelona und
Paris vom 1.\ bis zum 10.\ Februar 2027 und bestehen die abschließende
Prüfung der modellierten Nebenbedingungen. Die Bezeichnung \glqq Optimum\grqq{}
meint dabei den jeweiligen Bestwert der gefundenen Alternativen, keinen
Nachweis globaler Optimalität.

\begin{table}[H]
  \centering
  \small
  \begin{tabular}{@{}lrrr@{}}
    \toprule
    Auswahl & Kosten in Euro & Diskomfort & CO$_2$ in kg \\
    \midrule
    Kostenoptimum     & 228,51 & 496,7 & 321,9 \\
    Komfortoptimum    & 685,17 & 144,0 & 251,5 \\
    Ökologisches Optimum & 636,07 & 578,9 & 110,1 \\
    Kompromiss        & 361,37 & 316,6 & 252,8 \\
    \bottomrule
  \end{tabular}
  \caption{Empfehlungen des aktuellen Laufs; kleinere Zielwerte sind besser}
  \label{tab:eval_empfehlungen}
\end{table}

Gegenüber der ersten sichtbaren Ausgabe nach fünf Sekunden sinkt der
Kostenbestwert von 468,58 auf 228,51 Euro um 51,2\,\%, der geringste
Diskomfort von 311,9 auf 144,0 um 53,8\,\% und der Emissionsbestwert von
149,7 auf 110,1 kg um 26,5\,\%. Diese spaltenweisen Minima gehören nicht
notwendig zum selben Plan. Da Generation~0 keine Werte ausweist, ist dies
kein Vergleich mit der reinen Zufallsstartpopulation oder einer separaten
Zufallssuche. Frühe Fortschritte überwiegen, doch noch bei Generation~168
fällt der Preis von 267,38 auf 231,60 Euro; der letzte sichtbare
Komfortgewinn folgt erst bei Generation~341. Ein früher Plateauabbruch
hätte daher Verbesserungen verpasst. Unveränderte Einzelminima belegen zudem
keinen Stillstand der gesamten Paretofront.

Die erwarteten Zielkonflikte sind deutlich: Das Kostenoptimum nutzt drei
Flüge und einen Nachtbus, erreicht aber erheblich höheren Diskomfort und
mehr Emissionen als das fast dreimal so teure Komfortoptimum. Dieses
kombiniert drei Bahnfahrten mit einem Flug. Kurze Flugzeiten sind nicht
automatisch komfortabel: Beim günstigsten Plan werden aus 1,5 Stunden
Flugzeit nach London einschließlich Flughafenwegen und Abfertigungsreserven
8,5 Stunden. Der Komfortplan benötigt auf dieser Etappe mit der Bahn
einschließlich Reserven rund 7,3 Stunden und vermeidet späte Ankünfte.
Die fehlenden Zubringerkosten können den Preisvorteil der Flüge allerdings
überzeichnen.

Das ökologische Optimum verzichtet vollständig auf Flüge, nicht jedoch
auf Busse: Es nutzt drei Bahnverbindungen und einen Bus. Seine Emissionen
liegen um 65,8\,\% unter denen des günstigsten Plans, also bei etwa einem
Drittel. Nachtverbindungen verlagern einen Teil der Fahrt aus den
Tagesstunden, erhöhen jedoch die Belastung: Mit 578,9 Punkten ist dies
die unkomfortabelste Empfehlung. Die aufsummierte nutzbare Tageszeit
beträgt 81,03 gegenüber 76,41 Stunden beim Kostenoptimum, ein moderater
Vorteil von 4,62 Stunden. Diese verfügbaren Zeitfenster sind keine
garantierten Besichtigungszeiten.

Der Kompromiss kostet 132,86 Euro mehr als die günstigste Reise, reduziert
dafür den Diskomfort um 36,3\,\% und die Emissionen um 21,5\,\%.
Mit zwei Flügen, einer Bus- und einer Bahnfahrt vermittelt er zwischen
Preis und Belastung, erreicht aber mit 70,29 Stunden weniger nutzbare
Tageszeit. Die Auswahl hängt von Gewichten und Skalen ab; sie ist kein
objektiv bester Plan für jede reisende Person.

Frühere Entwicklungsversuche motivierten Zwischenspeicher, ein historisches
Pareto-Archiv und transportbezogene Vielfaltssicherung; zusätzliche
Suchoperatoren zeigten dagegen nicht durchgehend Vorteile.
Geänderte Reisedaten, Restriktionen und Emissionsmodelle verhindern einen
direkten Vergleich historischer Bestwerte. Zur Beurteilung von Suchgüte
und Hyperparametern sind Wiederholungen mit mehreren Zufallsstartwerten
unter identischen Bedingungen erforderlich. Die Modellzulässigkeit
garantiert zudem keine Buchbarkeit; Schätzpreise und Emissionsannahmen
begrenzen die Übertragbarkeit.
